\subsection{Dataset and Data Processing}
We use the Oxford-IIIT Pet Dataset for pixel-level segmentation. The dataset contains 37 categories
of cats and dogs, with approximately 200 images per class. Following the official dataset split as we have mentioned in Task1.

For semantic segmentation, the original trimap annotations are converted into three classes:
\begin{itemize}
    \item background,
    \item pet foreground,
    \item boundary region.
\end{itemize}

The mask labels are remapped to class indices $\{0,1,2\}$ for multi-class segmentation training.

Meanwhile, all images and masks are resized to a fixed spatial resolution of $128 \times 128$. The RGB images are converted into tensors and normalized using the ImageNet mean and standard deviation:
\[
\mu = (0.485, 0.456, 0.406), \quad
\sigma = (0.229, 0.224, 0.225).
\]

For segmentation masks, nearest-neighbor interpolation is used during resizing in order to preserve discrete class labels.

The processed dataset is loaded using the PyTorch \texttt{DataLoader} with mini-batch training. The training data are shuffled at every epoch to improve training randomness.